% ID: SEL-D0012-BA-E01
% ESTADO: cotejada. Segunda pasada visual completa y borrador compilado revisados.
% Fuente: Selectividad/matematicas_aplicadas/EXTRA_RESERVA_B_Modelo-6-Enunciados (Mat-Ap-2021).pdf
% SHA-256: 1752ec2f46b4ef2cc5d1c3346d69b79c2cc8a7aa12f6a0cd8f0d994b5bcda157
% Página: 1; bloque A; ejercicio original 1.

\textbf{(2.5 puntos)} La Agencia Espacial Europea contará con un presupuesto
de 2.4 millones de euros para financiar misiones sobre Observación de la Tierra
y para financiar programas de Transporte Espacial. Cada misión supone una
inversión de 200\,000 euros y cada programa, 100\,000 euros. Teniendo en cuenta
que en la decisión final deben superarse los 2 millones de euros de inversión
y el número de misiones debe ser al menos 4, pero no más de la mitad del número
de programas, \textquestiondown{}cuántas misiones y cuántos programas se deben
llevar a cabo para obtener el máximo de la función $F(x,y)=0.6x+0.4y$,
con $x$ misiones e $y$ programas?
